% Portada de J3.1: un cronómetro y barras de tiempo de tres versiones de una función.
\documentclass[tikz,border=0pt]{standalone}
\usepackage{resonancia}
\begin{document}
\begin{tikzpicture}
  \clip[rounded corners=6pt] (0,0) rectangle (12,3.6);
  \shade[top color=rNaranja!14, bottom color=white] (0,0) rectangle (12,3.6);
  % cronómetro
  \draw[trazo, color=rTinta, fill=white, line width=1.6pt] (2.2,1.9) circle (1.2);
  \fill[rTinta] (2.0,3.2) rectangle (2.4,3.45);
  \foreach \a in {0,30,...,330} \draw[color=rGris, line width=0.8pt] (2.2,1.9) ++(\a:1.0) -- ++(\a:0.14);
  \draw[color=rNaranja, line width=1.6pt] (2.2,1.9) -- ++(50:0.95);
  \fill[rTinta] (2.2,1.9) circle (0.07);
  % barras de tiempo
  \node[etiqueta, anchor=west] at (4.6,3.05) {\small versión ingenua};
  \fill[rNaranja!75] (7.0,2.9) rectangle (11.4,3.2);
  \node[etiqueta, anchor=west] at (4.6,2.15) {\small con vistas};
  \fill[rAzul!75] (7.0,2.0) rectangle (9.2,2.3);
  \node[etiqueta, anchor=west] at (4.6,1.25) {\small sin asignaciones};
  \fill[rVerde!75] (7.0,1.1) rectangle (7.9,1.4);
  \node[etiqueta, color=rGris] at (8.6,0.45) {\small medir, cambiar una cosa y volver a medir};
\end{tikzpicture}
\end{document}
